\section{Background and Related Work}
\label{sec:related}

\subsection{Approximate Nearest Neighbor Search}
\label{sec:ann}

Given a dataset $\mathcal{D} \subset \mathbb{R}^d$ of $n$ points, the nearest neighbor search problem asks to find the point $p^* \in \mathcal{D}$ closest to a query point $q$ under a distance metric $d(\cdot, \cdot)$.
Exact nearest neighbor search requires $\Omega(nd)$ time in the worst case, which becomes prohibitive for large-scale, high-dimensional datasets common in modern applications such as retrieval-augmented generation (RAG), recommendation systems, and embedding-based search.

The approximate nearest neighbor (ANN) problem relaxes this requirement by allowing a bounded error.

\begin{definition}[$(1+\epsilon)$-Approximate Nearest Neighbor]
\label{def:epsilon-ann}
Given a metric space $(X, d)$, a dataset $\mathcal{D} \subset X$, and a query point $q \in X$, let $p^* \in \mathcal{D}$ be the true nearest neighbor of $q$ such that $d(q, p^*) \le d(q, p)$ for all $p \in \mathcal{D}$. For any $\epsilon > 0$, a point $\hat{p} \in \mathcal{D}$ is a $(1+\epsilon)$-approximate nearest neighbor of $q$ if:
\begin{equation}
d(q, \hat{p}) \le (1 + \epsilon) \, d(q, p^*).
\end{equation}
\end{definition}

The ANN problem has been approached through several paradigms.
\textbf{Locality-Sensitive Hashing (LSH)}~\cite{indyk1998,datar2004} maps nearby points to the same hash bucket with high probability, enabling sub-linear query time with theoretical guarantees but requiring substantial space.
\textbf{Space-partitioning trees} such as KD-trees~\cite{bentley1975} and random projection trees (RP-trees)~\cite{dasgupta2008} recursively partition the space, offering efficient low-dimensional search but degrading in high dimensions due to the curse of dimensionality.
\textbf{Product quantization (PQ)}~\cite{jegou2011pq} compresses vectors into compact codes by decomposing the space into Cartesian products of subspaces, enabling fast approximate distance computation with reduced memory.
\textbf{Graph-based methods}~\cite{malkov2020,dong2011} construct proximity graphs where nodes are data points and edges connect approximate neighbors, supporting greedy traversal for search.
Among these, graph-based approaches have emerged as the dominant paradigm, offering superior recall-throughput trade-offs on standard benchmarks~\cite{annbenchmarks}.

\subsection{Construction Recall vs.\ Search Recall}
\label{sec:construction_recall}

In graph-based ANN, it is important to distinguish between \emph{construction recall} and \emph{search recall}, as they measure fundamentally different quantities.

\textbf{Construction recall} measures the fraction of true $k$-nearest neighbors that appear in each node's neighbor list in the constructed graph:
\begin{equation}
    \text{Recall}_{\text{constr}} = \frac{1}{n} \sum_{i=1}^{n} \frac{|N_G(i) \cap N^*(i)|}{k}
    \label{eq:construction_recall}
\end{equation}
where $N_G(i)$ is the set of $k$ neighbors of point $i$ in the constructed graph and $N^*(i)$ is the set of its true $k$-nearest neighbors computed via brute force.
This metric directly evaluates the quality of the graph construction algorithm---a construction recall of 0.6 means that, on average, 60\% of each node's true nearest neighbors are present in its neighbor list.

\textbf{Search recall}, by contrast, measures the fraction of true $k$-nearest neighbors returned by a search algorithm (such as greedy beam search) that operates on the constructed graph at query time:
\begin{equation}
    \text{Recall}_{\text{search}} = \frac{1}{|\mathcal{Q}|} \sum_{q \in \mathcal{Q}} \frac{|S(q) \cap N^*(q)|}{k}
    \label{eq:search_recall}
\end{equation}
where $\mathcal{Q}$ is the set of queries, $S(q)$ is the set of $k$ results returned by the search algorithm for query $q$, and $N^*(q)$ is the set of true $k$-nearest neighbors of $q$ in the dataset.

A graph with moderate construction recall (e.g., 0.60) can still achieve high search recall (e.g., 0.90+) because the search algorithm explores beyond immediate neighbors through multi-hop greedy traversal, effectively compensating for missing edges in the local neighborhood.
In this work, we primarily report construction recall, as it directly reflects the quality of the graph produced by NN-Descent and is the metric most affected by our filtering approach.

\subsection{Queries Per Second (QPS)}
\label{sec:qps}

Queries per second (QPS) is the standard throughput metric for evaluating ANN search systems.
It measures the number of queries that the system can answer per unit time:
\begin{equation}
    \text{QPS} = \frac{|\mathcal{Q}|}{T_{\text{search}}}
\end{equation}
where $|\mathcal{Q}|$ is the number of queries processed and $T_{\text{search}}$ is the total wall-clock time for all queries.

In practice, there is a fundamental trade-off between QPS and recall: exploring more of the graph at query time increases the probability of finding true neighbors (higher recall) but requires more distance computations (lower QPS).
The standard way to evaluate ANN systems is through Pareto plots of recall versus QPS, where a system dominates another if it achieves higher recall at the same QPS, or higher QPS at the same recall~\cite{annbenchmarks}.
Parameters such as the search beam width or exploration factor control this trade-off and are typically swept to generate the full Pareto curve.

\subsection{NN-Descent}
\label{sec:nndescent}

NN-Descent, proposed by Dong et al.~\cite{dong2011}, is a graph-based algorithm for constructing approximate $k$-nearest neighbor ($k$-NN) graphs.
The core intuition is simple: \emph{a neighbor of a neighbor is likely also a neighbor}.
Starting from a random graph, the algorithm iteratively refines the neighbor lists by comparing each point's neighbors with each other, exploiting the transitivity of proximity.

\subsubsection{Algorithm Overview}

The algorithm proceeds as follows:
\begin{enumerate}
    \item \textbf{Initialization:} Each point is assigned $k$ random neighbors, forming an initial (poor) approximation of the $k$-NN graph.
    \item \textbf{Local Join:} For each point $u$, all pairs $(v_1, v_2)$ drawn from $u$'s neighbor list are compared. If $d(v_1, v_2)$ is smaller than the distance to $v_1$'s or $v_2$'s current farthest neighbor, the respective neighbor list is updated.
    \item \textbf{New/Old Partitioning:} Neighbors are tagged as \textsc{New} (added in the current or previous iteration) or \textsc{Old} (unchanged). Only pairs involving at least one \textsc{New} neighbor are evaluated, as \textsc{Old}--\textsc{Old} pairs have already been compared. This avoids redundant distance computations.
    \item \textbf{Sampling:} To limit the number of candidate pairs, each point samples at most $\rho k$ new neighbors and $\rho k$ old neighbors for the local join, where $\rho \in (0, 1]$ is the sampling rate.
    \item \textbf{Convergence:} The algorithm terminates when the fraction of neighbor list updates per iteration falls below a threshold $\delta$ (typically $\delta = 0.001$), indicating that the graph has stabilized.
\end{enumerate}

Dong et al.\ report an empirical complexity of $O(n^{1.14})$ distance computations for constructing the $k$-NN graph, making NN-Descent substantially faster than the brute-force $O(n^2 d)$ approach.
The algorithm's simplicity and embarrassingly parallel local-join structure make it well-suited for parallelization, which has led to its adoption as the construction backbone in systems such as CAGRA~\cite{ootomo2024} for GPU-accelerated graph building.

Despite its efficiency, the dominant cost of NN-Descent is the \emph{exact distance computation} in the local join.
As the graph converges, the vast majority of candidate pairs do not lead to graph updates, yet each still incurs the full $O(d)$ cost.
This waste problem, which we quantify in Section~\ref{sec:waste}, is the primary motivation for our work.

\subsection{Improvements on NN-Descent}
\label{sec:nndescent_improvements}

As Figure~\ref{fig:waste} illustrates, NN-Descent with random initialization starts with a graph of nearly zero construction recall and requires five to six iterations just to reach the $\approx 0.3$ recall level that tree-based initialization attains directly.
Starting from such a low-quality graph has several undesirable consequences.
First, virtually every candidate pair examined in the early iterations can trigger a graph update, which in the worst case allows the number of updates to grow on the order of $k^2 n$, and for variants that retain more than $k$ candidates per point this growth is even larger.
Second, the quality of the final graph is sensitive to the random seed used for initialization: different seeds can lead to noticeably different convergence trajectories and final recall.
Tree-based space-partitioning structures, although not competitive as standalone indexes in high dimensions, are nevertheless very effective at separating \emph{clearly distant} vectors at low cost.
This makes them well-suited for initialization: they produce a starting graph with non-trivial recall, reduce the number of refinement iterations required, and largely eliminate the seed-dependence of random initialization.
Additionally, the original NN-Descent paper~\cite{dong2011} focuses solely on graph construction and does not specify a search procedure; the tree structures built for initialization can be naturally reused at query time to provide entry points for graph-based search, closing this gap.
Two systems, PyNNDescent and EFANNA, build on these observations using different tree structures and have become the standard reference implementations of NN-Descent in practice.

\subsubsection{PyNNDescent}
\label{sec:pynndescent}

PyNNDescent~\cite{pynndescent}, developed by McInnes as part of the UMAP ecosystem, initializes the $k$-NN graph using multiple \textbf{random projection trees (RP-trees)}~\cite{dasgupta2008}.
An RP-tree recursively partitions the data along random hyperplanes, so points falling into the same leaf are likely to be close in the original space.
Using the leaf co-occurrences from $L$ trees (typically $L = 10$ with leaf size $\approx 5k$) yields an initial graph with roughly 30\% construction recall, which NN-Descent then refines.
At query time, the same RP-trees are reused to identify candidate entry points, after which a greedy beam search is performed on the graph.

\subsubsection{EFANNA}
\label{sec:efanna}

EFANNA, proposed by Fu and Cai~\cite{fu2016}, follows the same hybrid recipe but uses multiple randomized \textbf{truncated KD-trees} instead of RP-trees.
A truncated KD-tree terminates subdivision once a leaf contains $K$ points (rather than a single point), which makes tree construction substantially faster than classical KD-trees.
The initial graph is constructed via a hierarchical divide-and-conquer procedure over these trees, and is then refined with standard NN-Descent iterations.
Like PyNNDescent, EFANNA also provides a search algorithm that uses the KD-trees to generate initial candidates, followed by NN-expansion on the graph.

A particularly relevant empirical finding from the EFANNA paper is that \emph{moderate losses in construction recall do not translate into proportional losses in search recall}.
Specifically, Fu and Cai report that on SIFT~1M an EFANNA index built on a graph with only $60\%$ construction recall still achieves search recall comparable to using the exact ground-truth $k$-NN graph, and substantially better than tree-only indexes such as FLANN~\cite{muja2014}.
The reason is that the ``missing'' neighbors in an approximate graph are not random points: on SIFT~1M, $98.9\%$ of the entries in a $60\%$-accurate $10$-NN graph constructed by EFANNA still lie within the true $100$ nearest neighbors.
In other words, the errors are ``near misses'' rather than outliers, and greedy graph search can easily navigate past them through multi-hop traversal.
This observation is important for our work: it suggests that small reductions in construction recall induced by a filter during graph building can be tolerated, since they have a disproportionately small effect on downstream search quality.

Despite these improvements, neither PyNNDescent nor EFANNA addresses the core inefficiency of the NN-Descent iterations themselves---namely, the high volume of exact distance computations that do not update the graph.
Neither system applies any form of distance filtering inside the local-join loop, which is the gap our work targets.

\subsection{Filtering in Graph-Based ANN}
\label{sec:filtering}

Recent work has explored the idea of using cheap approximate distance estimates to \emph{filter} candidate pairs before computing their exact distances, thereby reducing the computational cost of graph construction or search.
We discuss two approaches that directly motivate our work.

\subsubsection{LSH-APG}
\label{sec:lsh_apg}

LSH-APG (Locality-Sensitive Hashing for Approximate Proximity Graph construction), proposed by Zhao et al.~\cite{zhao2023}, integrates locality-sensitive hashing into proximity graph construction.
Their key contribution relevant to our work is a \textbf{projection-based distance filter} that uses random projections and chi-squared concentration bounds to cheaply identify candidate pairs that are unlikely to be neighbors.

The filter works as follows.
Given $m$ random Gaussian vectors $\mathbf{a}_1, \ldots, \mathbf{a}_m \sim \mathcal{N}(0, I_d)$, each data point $x \in \mathbb{R}^d$ is projected to a lower-dimensional representation:
\begin{equation}
    P(x) = \big(\mathbf{a}_1 \cdot x, \;\; \mathbf{a}_2 \cdot x, \;\; \ldots, \;\; \mathbf{a}_m \cdot x\big) \in \mathbb{R}^m
    \label{eq:projection}
\end{equation}
These projections are computed once during preprocessing at a cost of $O(n \cdot m \cdot d)$.

The theoretical foundation rests on the \textbf{concentration of projected distances}.
For any two points $u, v \in \mathbb{R}^d$ with Euclidean distance $\|u - v\|_2$, the squared projected distance satisfies:
\begin{equation}
    \frac{\|P(u) - P(v)\|^2}{\|u - v\|_2^2} \sim \chi^2(m)
    \label{eq:chi2}
\end{equation}
where $\chi^2(m)$ denotes the chi-squared distribution with $m$ degrees of freedom.
This follows from the fact that each component $\mathbf{a}_i \cdot (u - v)$ is an independent Gaussian with variance $\|u - v\|_2^2$, so the sum of their squares (after normalization) follows a chi-squared distribution.

By the chi-squared concentration inequality, for a confidence level $p_\tau \in (0, 1)$:
\begin{equation}
    \Pr\!\Big[\|P(u) - P(v)\|^2 \leq \chi^2_{p_\tau}(m) \cdot \|u - v\|_2^2\Big] \geq p_\tau
    \label{eq:concentration}
\end{equation}
where $\chi^2_{p_\tau}(m)$ is the $p_\tau$-quantile of the chi-squared distribution with $m$ degrees of freedom.
In words: with probability at least $p_\tau$, the squared projected distance is no larger than $\chi^2_{p_\tau}(m)$ times the true squared distance.

The contrapositive gives the filtering criterion: if the projected distance exceeds $\chi^2_{p_\tau}(m) \cdot r^2$ for a reference distance $r$, then with confidence $p_\tau$, the true distance exceeds $r$, and the pair can be safely skipped.
LSH-APG applies this filter during HNSW-style graph construction, achieving significant reductions in distance computations during the neighbor candidate evaluation phase.

Our work adapts this projection-based filtering idea specifically to the \emph{local-join} structure of NN-Descent, where the reference distance $r$ is naturally defined as the current farthest neighbor distance $d_k(u)$ for each endpoint, and the filter is applied with a conservative conjunctive (AND) condition requiring both endpoints to agree that the pair is far.

\subsubsection{FINGER}
\label{sec:finger}

FINGER (Fast Inference for Graph-based approximate nearest neighbor search using projections for the Estimation of distanRes), proposed by Chen et al.~\cite{chen2023finger}, takes a complementary approach: rather than filtering during graph \emph{construction}, it accelerates graph \emph{search} (query time) on an already-built HNSW index.

The key idea is to precompute a low-rank approximation of the dataset and use it to cheaply estimate distances during the greedy beam search.
Specifically, FINGER computes a rank-$r$ approximation of the data matrix using randomized SVD:
\begin{equation}
    X \approx U_r \Sigma_r V_r^\top
\end{equation}
where $U_r \in \mathbb{R}^{n \times r}$, $\Sigma_r \in \mathbb{R}^{r \times r}$, and $V_r \in \mathbb{R}^{d \times r}$.
For each data point, the projected representation $\tilde{x}_i = V_r^\top x_i \in \mathbb{R}^r$ is stored.

During search, for a candidate pair $(q, x_i)$, FINGER first computes the projected distance $\|\tilde{q} - \tilde{x}_i\|^2$ in $O(r)$ time.
Using concentration bounds on the residual (the component orthogonal to the projection subspace), FINGER derives a \emph{lower bound} on the true distance:
\begin{equation}
    \|q - x_i\|^2 \geq \|\tilde{q} - \tilde{x}_i\|^2 + \text{LB}_{\text{residual}}
\end{equation}
If this lower bound exceeds the current search radius, the candidate is pruned without computing the full $O(d)$ distance.

FINGER reports 20--60\% speedups on high-dimensional datasets while maintaining the same recall as standard HNSW search.
The approach is complementary to ours: FINGER reduces distance computations during \emph{search} on a fixed graph, while our work reduces distance computations during \emph{graph construction} via NN-Descent.
In principle, both techniques could be applied simultaneously---our filter during construction and FINGER during search---to accelerate the full pipeline.
